\documentclass[12pt, letterpaper]{article}
\usepackage[margin=1.2in]{geometry}
\usepackage{ebgaramond}
\usepackage[T1]{fontenc}
\usepackage{microtype}
\usepackage{setspace}
\usepackage{parskip}
\usepackage{titling}
\usepackage{fancyhdr}

\setlength{\parindent}{2em}
\onehalfspacing

\pagestyle{fancy}
\fancyhf{}
\fancyhead[L]{\small\textit{The Thermometer}}
\fancyhead[R]{\small\textit{NASA Mission Series v1.12}}
\fancyfoot[C]{\small\thepage}
\renewcommand{\headrulewidth}{0.4pt}

\title{\Huge\textbf{The Thermometer}\\[0.5em]\large\textit{A short story from Mission 1.12 --- Explorer 7}}
\author{NASA Mission Series\\[0.3em]\small Miles Tiberius Foxglove --- Mukilteo, WA}
\date{March 30, 2026}

\begin{document}
\maketitle
\thispagestyle{empty}

\vspace{1em}
\noindent\rule{\textwidth}{0.4pt}
\vspace{1em}

\noindent The data came down as numbers. Not pictures, not images, not anything a person could look at and understand without translation. Just numbers: voltage readings from two small hemispheres bolted to the skin of a satellite that weighed less than a large dog, orbiting the Earth at 28,000 kilometers per hour in an ellipse that ranged from 573 to 1,073 kilometers above the surface. The numbers arrived as modulated tones on a radio frequency of 108 megahertz, received by ground stations, recorded on strip chart paper, and delivered to a professor in Wisconsin who had designed the instruments that produced them.

Verner Suomi opened the first data package on a Tuesday morning in his office at the University of Wisconsin-Madison. The strip charts were long, narrow rolls of paper with a wiggly line traced by a pen arm that deflected in response to voltage changes. One chart was labeled BLACK HEMISPHERE and the other WHITE HEMISPHERE. The charts covered six orbits --- approximately ten hours of data from the first day of Explorer 7's operation, October 13, 1959.

He unrolled the black hemisphere chart first and weighted the ends with coffee cups. The trace rose and fell with a period of approximately 101 minutes --- the orbital period. When the satellite was on the sunlit side of the Earth, the black hemisphere absorbed everything: direct sunlight, sunlight reflected from Earth's clouds and surface, and thermal infrared radiation emitted by the Earth. The trace climbed. When the satellite crossed the terminator into Earth's shadow, the solar components disappeared and only the thermal emission remained. The trace dropped. The pattern was regular, predictable, and exactly what the physics said it should be.

Then he unrolled the white hemisphere chart and placed it alongside the black.

\begin{center}
$\bullet\quad\bullet\quad\bullet$
\end{center}

The white hemisphere was the trick. The insight that made everything work. Suomi and his colleague Robert Parent had coated one hemisphere with a surface that absorbed radiation of all wavelengths --- the black sensor --- and the other with a surface that reflected visible light but absorbed infrared --- the white sensor. When sunlight hit the white hemisphere, most of it bounced off. When infrared radiation from the Earth hit the same surface, it was absorbed. The white hemisphere was blind to sunshine but sensitive to heat.

This meant that the difference between the two traces was the solar component. The black hemisphere saw everything; the white hemisphere saw only the thermal part. Subtract white from black and what remained was the contribution of visible light --- direct solar radiation plus the sunlight reflected from Earth's surface and clouds. This was the albedo measurement. This was the number that determined how much of the sun's energy the Earth kept and how much it threw back.

Suomi traced his finger along the white hemisphere chart. It rose and fell too, but with a smaller amplitude. The thermal emission from the Earth was nearly constant --- the planet radiated at roughly the same rate whether the satellite was over ocean or land, day side or night side. The temperature of the outgoing radiation was approximately 255 Kelvin --- minus eighteen degrees Celsius --- the effective radiating temperature of the Earth as seen from space.

He already knew this number from theory. The Stefan-Boltzmann law predicted it. If you took the solar constant --- 1,361 watts per square meter, the amount of energy the sun delivered to a surface facing it at Earth's distance --- and multiplied by the fraction that Earth absorbed after reflecting 30 percent, and divided by four to account for the ratio of cross section to surface area, you got approximately 239 watts per square meter. Set that equal to $\sigma T^4$ and solve for $T$, and you got 255 Kelvin. Minus eighteen Celsius. Twenty degrees below freezing.

But the actual surface of the Earth was not minus eighteen. It was plus fifteen. Thirty-three degrees warmer than the bare radiative equilibrium.

\begin{center}
$\bullet\quad\bullet\quad\bullet$
\end{center}

The thirty-three degrees was the greenhouse effect, and Suomi was looking at it. Not inferring it from surface measurements. Not calculating it from atmospheric composition. Looking at it, from space, for the first time. The black hemisphere told him how much energy was arriving. The white hemisphere told him how much was leaving as thermal radiation. The difference between what came in and what went out was the energy trapped by the atmosphere --- by water vapor, carbon dioxide, methane, nitrous oxide, and the other gases that absorbed outgoing infrared and re-emitted it back toward the surface.

He picked up a pencil and began to annotate the charts. At the point where the satellite crossed from the day side to the night side, he wrote ``terminator.'' At the peak of the black hemisphere trace, he wrote the voltage and, next to it, the calibrated flux in watts per square meter. At the bottom of the white hemisphere trace during eclipse, he wrote the thermal emission alone, uncontaminated by any solar component. He drew a bracket between the two traces on the sunlit side and wrote: ``this is what the atmosphere keeps.''

The strip charts were the first planetary thermometer reading taken from outside the patient.

\begin{center}
$\bullet\quad\bullet\quad\bullet$
\end{center}

Rudolf Virchow had understood this principle a century earlier, though he applied it to a different kind of body. In 1858 --- one year before Explorer 7 launched --- Virchow published \textit{Cellular Pathology}, the work that established the principle that disease in the whole organism could be diagnosed by examining individual cells. You did not need to understand the entire body to know it was sick. You needed to look at the right cells, under the right conditions, and the cells would tell you. The part revealed the condition of the whole.

Suomi's bolometers were cellular pathology for the planet. Two hemispheres, each 2.5 centimeters in diameter. Two cells extracted from the body of the Earth's radiation field. One that absorbed everything, one that was selective. The difference between them diagnosed the energy imbalance as precisely as a blood test diagnosed anemia. The instrument weighed grams. The satellite weighed 41.5 kilograms. The patient weighed six trillion trillion kilograms. And the diagnosis was clear: the patient was retaining more energy than it was releasing. Not much --- less than one watt per square meter, on average. But the fever was real.

Virchow had also said: ``The physician is the advocate of the poor.'' He meant that medicine was not merely a technical practice but a social responsibility --- that the doctor's obligation extended beyond diagnosis to action, beyond understanding to intervention. Suomi would spend the next three decades advocating for the same principle applied to climate. The measurement was not enough. The measurement demanded a response.

\begin{center}
$\bullet\quad\bullet\quad\bullet$
\end{center}

In the forests of the Pacific Northwest, in the wet shade of old-growth Sitka spruce and western hemlock, \textit{Usnea longissima} hangs in pale curtains from the branches. Methuselah's beard lichen --- the longest lichen in the world, strands reaching three meters or more, draped like grey-green hair from the canopy. It grows slowly, a few centimeters per year at most. Some colonies are centuries old. It absorbs moisture and nutrients directly from the air, having no roots, no vascular system, no way to buffer itself against changes in atmospheric composition. It is a sponge for whatever the air contains.

This makes it a thermometer. Not a thermometer that measures temperature --- a thermometer that measures change. When the air quality shifts --- when sulfur dioxide increases, when ozone levels change, when particulate matter rises, when the moisture balance tips --- \textit{Usnea longissima} responds before almost any other organism in the forest. Its growth rate changes. Its color shifts from sage green to grey. Its strands thin and break. In areas where air quality has declined, the lichen is simply gone, its absence a diagnostic marker as clear as any cell under Virchow's microscope.

The lichen does not know it is a sensor. It is not designed to measure anything. It simply exists at the interface between atmosphere and organism, absorbing what is there, growing when conditions are right, dying when they are not. It is the forest's bolometer: a body that absorbs the environment's radiation --- chemical, physical, thermal --- and whose temperature reveals the balance.

\begin{center}
$\bullet\quad\bullet\quad\bullet$
\end{center}

Suomi worked through the data for three days. He built a table: orbit number, sun angle, black hemisphere flux, white hemisphere flux, difference, implied albedo, implied thermal emission, implied energy imbalance. The numbers told a story that the strip charts had only sketched. Over the sunlit hemisphere, the Earth reflected approximately 30 percent of incoming sunlight and absorbed 70 percent. Over oceans, the reflection was lower --- dark water absorbed more. Over clouds, it was higher --- white surfaces reflected more. Over deserts, over forests, over ice --- each surface type had its own albedo, and Explorer 7's orbit carried the bolometers over all of them.

The thermal emission was more uniform. The Earth radiated at roughly the same rate everywhere, because the atmosphere mixed and redistributed the energy before it escaped to space. Tropical oceans absorbed much more solar energy than polar ice, but the atmosphere carried the excess poleward in the form of weather systems, wind patterns, and ocean currents. The thermal emission was the equilibrated output --- the planet's temperature as averaged by its own fluid dynamics.

What Suomi saw in the data was the machine. The planet as a heat engine: solar input at the equator, thermal output at the poles, weather as the working fluid, albedo as the throttle, greenhouse gases as the insulation. The bolometers were the gauges. And the gauges showed that the engine was not quite in balance --- that the insulation was trapping slightly more than the output could release --- that the temperature was drifting, slowly, upward.

He could not prove this from Explorer 7 alone. The data covered weeks, not decades. The precision of the bolometers was limited. The orbital coverage was sparse --- one satellite, one inclination, one altitude range. But the principle was proven. The measurement could be made. The planet could be diagnosed from orbit, the way Virchow diagnosed a patient from a cell sample, the way \textit{Usnea longissima} diagnosed the air by growing in it or dying.

\begin{center}
$\bullet\quad\bullet\quad\bullet$
\end{center}

Suomi would go on to design better instruments. The Earth Radiation Budget Experiment, ERBE, flew on multiple satellites in the 1980s. The Clouds and the Earth's Radiant Energy System, CERES, began flying in 1997 and is still operating. Each generation of instruments improved the precision, expanded the coverage, extended the record. By the 2020s, the energy imbalance was measured at 0.87 watts per square meter --- less than one watt over every square meter of the planet's surface, but relentless, persistent, cumulative.

The lichen continued to disappear. In the lowland forests of the Pacific Northwest, \textit{Usnea longissima} became increasingly rare. Air quality changes, temperature increases, altered moisture patterns --- the compound effects of the energy imbalance that Suomi first measured in 1959 worked their way through the atmospheric chemistry and arrived at the surface of the lichen's cortex, where the photobiont and mycobiont lived in their ancient partnership, and the partnership faltered. The strands thinned. The color grayed. The curtains that had hung from old-growth branches for centuries shortened, broke, fell to the forest floor.

The lichen was the thermometer. The bolometer was the thermometer. Virchow's microscope was the thermometer. They were all the same instrument, applied at different scales to the same question: is the body in balance? Is the energy coming in equal to the energy going out? Is the system maintaining itself, or is it drifting?

The strip charts on Suomi's desk said: drifting. Not fast. Not dramatically. One watt per square meter is an imperceptible quantity at human scale --- you cannot feel it, cannot see it, cannot detect it without instruments purpose-built for the measurement. But integrated over the surface of a planet, over decades, over the thermal mass of oceans four kilometers deep, one watt per square meter is sufficient to raise the temperature of everything. The fever is real. The cells show it. The lichen shows it. The bolometer showed it first.

Verner Suomi unrolled the strip charts on a Tuesday morning and read the planet's temperature for the first time from outside. The thermometer worked. The patient was warm.

\end{document}
